% =========================================================================
% Figura: Sistema de coordenadas esféricas
% =========================================================================
% Archivo: Coordenadas_esfericas.tikz.tex
% Usada en: C1_00_introduccion.tex
% Convención del libro: theta = ángulo azimutal (plano xy), phi = ángulo polar
% =========================================================================

\begin{figure}[!ht]
    \centering
    \tdplotsetmaincoords{70}{110}
    \begin{tikzpicture}[tdplot_main_coords, >=Latex,
                        ejes/.style={->, thick, gris1},
                        univec/.style={->, very thick, azul},
                        posvec/.style={->, very thick, myred}]

        % --- CONFIGURACIÓN DE PARÁMETROS ---
        \pgfmathsetmacro{\phivec}{48.17}   % Ángulo polar [deg]
        \pgfmathsetmacro{\thetavec}{63.5}  % Ángulo azimutal [deg]
        \def\rlen{6}                       % Largo del vector posición [u]

        % --- SISTEMAS DE COORDENADAS ---
        \draw[ejes] (0,0,0) -- (6.5,0,0) node[pos=1.1] {$x$};
        \draw[ejes] (0,0,0) -- (0,6,0) node[pos=1.05] {$y$};
        \draw[ejes] (0,0,0) -- (0,0,5.5) node[pos=1.05] {$z$};

        % --- PUNTOS Y COORDENADAS ---
        \tdplotsetcoord{P}{\rlen}{\phivec}{\thetavec}

        % --- VECTORES Y FUERZAS ---
        % Vector posición
        \draw[posvec] (0,0,0) -- (P) node[midway, above left] {$\vec{r}$};
        % Vectores unitarios en el punto P
        \tdplotsetcoord{Er}{7}{\phivec}{\thetavec}
        \draw[univec] (0,0,0) -- (Er) node[pos=1.05] {$\hat{e}_r$};
        \tdplotsetcoord{Et}{1}{90}{90+\thetavec}
        \draw[univec] (Pxy) -- ($(Et) + (Pxy)$) node[pos=1.3] {$\hat{e}_\theta$};
        \tdplotsetcoord{Ep}{1}{90+\phivec}{\thetavec}
        \draw[univec] (P) -- ($(Ep) + (P)$) node[pos=1.3] {$\hat{e}_\phi$};

        % --- ANOTACIONES Y ETIQUETAS ---
        % Líneas de ayuda (proyecciones)
        \draw[thick] (0,0,0) -- (Pxy);
        \draw[mydashed, gris1!60] (P) -- (Pxy);
        \draw[mydashed, gris1!60] (Px) -- (Pxy) node[pos=-0.15] {$x$};
        \draw[mydashed, gris1!60] (Py) -- (Pxy) node[pos=-0.3] {$y$};
        \draw[mydashed, gris1!60] (Pz) -- (P) node[pos=-0.15] {$z$};
        % Punto P
        \node[fill=gris1, circle, inner sep=0.9pt, label={above right:$P$}] at (P) {};
        % Ángulos
        \tdplotdrawarc{(0,0,0)}{0.7}{0}{\thetavec}{below}{$\theta$}
        \tdplotsetthetaplanecoords{\thetavec}
        \tdplotdrawarc[tdplot_rotated_coords]{(0,0,0)}{0.5}{0}{\phivec}{}{}
        \node at (0,0.25,0.67) {$\phi$};

    \end{tikzpicture}
    \caption{Coordenadas esféricas y sus vectores unitarios en el punto $P$}
    \label{fig:c1_intro_coordenadas_esfericas}
\end{figure}
